\documentclass[12pt]{book}
\usepackage{amsmath}
\usepackage{amssymb}
\usepackage{geometry}
\geometry{margin=1in}

\title{Volume 12: Computation, Code, and Simulation Systems \\ Book 01: Computational Core}
\author{James C. Harvey}
\date{December 2025}

\begin{document}
\maketitle
\tableofcontents

\chapter{Code as Formalization}

\section*{Abstract}
The computational core of SDT is a formalization of the geometric equations. This book documents how
constants, operators, and solvers implement the master equation and state-space operators.

\section{Constants and Units}
The constant set used across implementations anchors the numerical representation of SDT geometry.

\section{Operator Implementation}
State operators are encoded to preserve geometric meaning: curvature, circulation, slip, and occlusion
are evaluated numerically but remain SDT-native quantities.

\chapter{State Vector Implementation}

\section*{Abstract}
The 28D state vector is implemented explicitly so that each component maps to a geometric quantity. This
chapter describes the numerical representation and its usage.

\section{Component Mapping}
Each $\xi$, $T$, $\Phi$, and $\epsilon$ component is stored with unit metadata to preserve dimensional
clarity during computation.

\chapter{Units and Dimensional Checks}

\section*{Abstract}
Dimensional checks are mandatory for every operator. This chapter documents the internal unit discipline.

\section{Unit Assertions}
Each function asserts dimensional validity:
\begin{equation}
[\dot{E}] = [P] [A] [\Gamma] [\kappa],
\end{equation}
ensuring the master equation is consistent in code.

\chapter{Testing and Validation Harness}

\section*{Abstract}
Code correctness is enforced through tests that reproduce canonical benchmark values.

\section{Harness Structure}
The harness loads constants, constructs state vectors, evaluates operators, and compares outputs to
reference artifacts.
\end{document}
